\textbf{COLUNI 2026 -- Questão 35}\par\medskip Em um de seus escritos, Mahatma Gandhi afirmou: “\textit{A não violência, em sua concepção dinâmica, significa sofrimento consciente. Não quer absolutamente dizer submissão humilde à vontade do malfeitor, mas um empenho, com todo o ânimo, contra o tirano.}”\par
Fonte: GANDHI, Mahatma. \textbf{Pensamentos de Mahatma Gandhi}. São Paulo: Martin Claret, 2002.

Considerando essa afirmação e os acontecimentos ligados à trajetória de Gandhi e ao processo de independência da Índia, assinale a afirmativa CORRETA:\begin{enumerate}
\item A Satyagraha, doutrina da desobediência civil defendida por Gandhi, consistia no uso moderado da força militar indiana contra os colonizadores britânicos.
\item A atuação de Gandhi no processo de independência da Índia foi precedida por um tempo na África do Sul, quando o líder indiano, mesmo convivendo com situações de racismo, manteve-se fora da luta social.
\item O contato de Gandhi com o racismo na África do Sul fortaleceu sua visão de que a opressão colonial não poderia ser combatida por meio da violência, mas sim pela firmeza moral e pela desobediência civil.
\item O processo de independência da Índia realizou o sonho de Gandhi de uma nação sem divisões religiosas, onde hindus e muçulmanos vivessem em harmonia sob o governo de um mesmo país.\end{enumerate}